\subsection{NN-G024: nonprimitive common cells give the exact 5932-pair optimum}
\paragraph{Statement.}
There is an effective associative rank-three dimension ring with
$5932$ cancelling pairs, of which $5208$ are in the bottom cell.
Moreover $5932$ is the exact minimum in the entire following
class: the closed rank-three table shape, a common cell polynomial
$\delta\in\mathbb Z[D]$ of arbitrary degree and coefficient content,
arbitrary-degree $W\in\mathbb Z[D]$, positive rational constant
ratios $G_i/(W\delta)$, and arbitrary rational auxiliary functions
$\Lambda_i\in\mathbb Q(D)$. Only the actual multiplication
coefficients are required to be integral, effective, coefficientwise
dominant over the absolute trace targets, and congruent to those
targets modulo two.

This class includes all common-cell constructions in NN-G019,
NN-G021 and NN-G023 and removes their monicity and cell-degree
restrictions. The new optimum is attained by a cell polynomial
of content three and a nonintegral auxiliary $\Lambda_2$.
The theorem does not cover unequal general cells, nonconstant
ratios, or other table shapes.

\paragraph{An explicit integral table with a rational auxiliary factor.}
Put $\delta_0=(D-1)(D-2)$ and choose
\[
 \begin{aligned}
 \delta&=3\delta_0,& W&=3D,\\
 G_1&=6D\delta_0,&G_2&=9D\delta_0,&S&=18D\delta_0^2,\\
 U_1&=D,&U_2&=3D,\\
 \Lambda_1&=D-2,&\Lambda_2&=(D^2-2D-4)/3,\\
 V_1&=2D^2-D+2,&V_2&=4D^2+D+2.
 \end{aligned}                                                   \tag{G024.1}
\]
The ratios are $(k_1,k_2)=(2/3,1)$. The identities
\[
 S=G_1G_2/W,\qquad
 U_i=(k_iW-D)/k_j,\qquad
 V_i=(k_iDW-\Lambda_i)/k_j\quad(j\ne i)
\]
hold exactly. The two matrix blocks and two scalar blocks in
NN-G011 therefore give an injective associative realization
of the closed table over $\mathbb Q(D)$: neither
$\Lambda_i(D^2-\Lambda_i)$ vanishes identically.

For effectivity and integrality it is the actual product
$\delta\Lambda_i$ that matters. In symmetric Laurent
coefficient lists indexed by nonnegative exponents, the base
coefficients of this table and their absolute trace targets are
\[
\begin{array}{c|l|l}
 &\text{dimension coefficient}&\text{absolute trace target}\\\hline
 \delta &(6,3,3)&(0,1,1)\\
 \delta\Lambda_1&(6,9,3,3)&(2,1,1,1)\\
 \delta\Lambda_2&(2,5,2,3,1)&(2,3,2,1,1)\\
 W&(6,3)&(2,1)\\
 U_1&(2,1)&(2,1)\\
 U_2&(6,3)&(2,1)\\
 V_1&(12,7,2)&(6,5)\\
 V_2&(28,17,4)&(6,3,2)\\
 G_1&(36,30,18,6)&(16,12,4)\\
 G_2&(54,45,27,9)&(14,11,5,1)\\
 S&(504,450,324,180,72,18)&(144,120,68,24,4)
\end{array}                                                       \tag{G024.2}
\]
Every left list is nonnegative, coefficientwise at least the
right list, and congruent to it modulo two. The closed table
uses these coefficients multiplied by nonnegative powers of
$D$. For signed trace cell coefficients use
$|D^\alpha f|\le D^\alpha|f|$ coefficientwise. These facts
prove every dimension-ring condition. In particular the
nonintegral $\Lambda_2$ is not itself a structure constant;
its product with the nonprimitive cell is integral.

At $D=4$ the bottom mass is
\[
 2592+20(144+216)+200\cdot12+20(4+12+30+70)=14512.
\]
The two cell masses sum to
$52\cdot18+14\cdot6(6+4)=1776$, so $T=16288$.
Subtracting the true absolute masses gives
\[
 (16288-4424)/2=5932,\qquad (14512-4096)/2=5208.
 \tag{G024.3}
\]
The independent integer-Laurent verifier
\texttt{verify\_dimension\_5932.py} reconstructs all $132$
positions and checks all $1000$ associativity triples. It
handles the actual integer products directly and does not
round or formally clear an unverified denominator.

\paragraph{Content normalization for the entire common-cell class.}
Write
\[
 \delta=c\delta_0,\quad W=gW_0,
\]
where $c,g>0$ are coefficient contents and $\delta_0,W_0$
are primitive symmetric integral Laurent polynomials.
They are also in $\mathbb Z[D]$, because symmetric integral
Laurent polynomials have an integral triangular expansion
in $q^j+q^{-j}$ and hence in $D=q+2+q^{-1}$.
The actual parity conditions imply that $c,g$ are odd,
\[
 \delta_0\equiv D(D-1),\qquad W_0\equiv D\pmod2.
\]
Both polynomials are effective. In particular
$d:=\delta_0(4)\ge4$ and $w:=W_0(4)\ge4$ are even.

Gauss's content lemma applied after Laurent shifts gives
positive integers $a_i$ such that
\[
 \begin{aligned}
 G_i&=a_iW_0\delta_0,&
 S&=(a_1a_2/g)W_0\delta_0^2,\\
 U_i&=g(a_iW_0-cD)/a_j,&
 V_i&=g(a_iDW_0-\Psi_i)/a_j,
 \end{aligned}                                                     \tag{G024.4}
\]
where $\Psi_i=c\Lambda_i$ and $j\ne i$.
Here $a_i=k_i g c$, $a_1$ is even, $a_2$ is odd, and
$g\mid a_1a_2$. Indeed the primitive products force the
integrality statements, and reduction modulo two of the
$G_i$ gives their different scalar parities.

Even though $\Lambda_i$ may be rational, $\Psi_i$ is in
$\mathbb Z[D]$. First (G024.4), with integral actual $V_i$,
puts $\Psi_i$ in $\mathbb Q[D]$. The actual within-cell
coefficient is $\delta\Lambda_i=\delta_0\Psi_i\in\mathbb Z[D]$.
Primitivity of $\delta_0$ and Gauss's lemma force
$\Psi_i\in\mathbb Z[D]$. This is the exact integrality
statement; it does not force division by the content $c$.
Modulo two these actual cell products also give
$\Psi_1\equiv D$ and $\Psi_2\equiv D^2$.

Set $p_i=\Psi_i(4)$. The exact mass expression is
\[
 \begin{split}
 T={}&(a_1a_2/g)wd^2+20wd(a_1+a_2)+200gw\\
 &+20\sum_i\{U_i(4)+V_i(4)\}+52cd+14d(p_1+p_2).
 \end{split}                                                       \tag{G024.5}
\]
The two effective cell products imply $dp_1\ge8$ and
$dp_2\ge16$. Also $U_i(4)\ge4$ and $V_i(4)\ge16$.
Thus the score bound of NN-G021 applies with an extra
nonnegative $52(c-1)d$ term. For any proposed improvement
on (G024.3) it gives the following generous complete bounds:
\[
 d\le30,\quad w\le36,\quad a_1+a_2\le52,\quad g\le23.
 \tag{G024.6}
\]
Furthermore the right side of (G024.5) is at least
$3024+208c$: use $a_1a_2/g\ge2$, $a_1+a_2\ge3$,
$d,w\ge4$, and the $1136$ lower bound for the last
cell-product term and the four $U,V$ terms. Since a
proposed improvement has $T<16288<18552$, the weaker
bound $c\le73$, with $c$ odd, is valid. None of these
bounds asserts that a Laurent support index is bounded.

\paragraph{A complete finite relaxation, retaining arbitrary support locations.}
Write
\[
 W_0=c_0+c_1(q+q^{-1})+\sum_{j\ge2}c_j(q^j+q^{-j}).
\]
Here $c_0\ge2$ is even, $c_1\ge1$ is odd, and $c_j$
is nonnegative even for $j\ge2$. The same outer-coefficient
argument as in (G023.2) gives $M\mid c_j$, where
\[
 M=\operatorname{lcm}\left(
 \frac{2a_2}{\gcd(ga_1,2a_2)},
 \frac{2a_1}{\gcd(ga_2,2a_1)}\right),\quad 4\mid M.
\]
Put $h=\sum_{j\ge2}c_j/M$, so $w=c_0+2c_1+2Mh$.
For $(b,t)=(c_0,2),(c_1,1)$ the necessary zero and
first coefficient conditions on $U_i$ are
\[
 g(a_i b-ct)/a_j\in t+2\mathbb N.                     \tag{G024.7}
\]
Primitivity requires $\gcd(c_0,c_1)=1$ for $h=0$ and
$\gcd(c_0,c_1,M)=1$ for $h>0$.

Let $a_j'=a_j/\gcd(g,a_j)$. Evaluating the $V_i$
congruences and domination gives the integer intervals
\[
 \begin{gathered}
 p_i\equiv4a_iw\pmod{2a_j'},\\
 \left\lceil X_i(4)/d\right\rceil\le p_i\le
 \left\lfloor(4ga_iw-16a_j)/g\right\rfloor,
 \quad(X_1(4),X_2(4))=(8,16).
 \end{gathered}                                                       \tag{G024.8}
\]
For fixed other data the coefficient of $p_i$ in (G024.5)
is $14d-20g/a_j$. Its minimum over (G024.8) is therefore
at the first allowed integer when this coefficient is
nonnegative and at the last otherwise. This is a lower
bound using necessary conditions, not an assertion that
such a polynomial $\Psi_i$ exists.

There is one additional necessary leading-coefficient
condition when $h=0$. Then $\deg W_0=1$. If $\deg\Psi_2>2$,
its highest Laurent coefficient would have opposite signs
in the effective polynomials $\delta_0\Psi_2$ and $V_2$.
The leading coefficient of $\delta_0$ is positive, so this
is impossible. Since $\Psi_2\equiv D^2\pmod2$, its
coefficient $L$ at $q^2$ must therefore be positive odd.
The coefficient of $V_2$ at $q^2$ is required to be an
even integer at least two. Consequently
\[
 L\ge1,\quad L\equiv a_2c_1-2a_1'\pmod{2a_1'},\quad
 L\le\left\lfloor(g a_2c_1-2a_1)/g\right\rfloor.
 \tag{G024.9}
\]
This test uses no cell-degree restriction.

For reproducibility the whole finite range is as follows:
even $2\le a_1\le50$, odd $1\le a_2\le51$ with sum
at most $52$; odd $1\le g\le23$ dividing $a_1a_2$;
odd $1\le c\le73$; even $4\le d\le30$;
even $2\le c_0\le36$; odd $1\le c_1\le17$;
and $0\le h\le(36-c_0-2c_1)/(2M)$.
Impose (G024.7), the primitive-content necessary tests,
$G_i(4)\ge48$, $S(4)\ge576$, and the endpoint lower
bound from (G024.8). The included integer-only script
\texttt{verify\_common\_cell\_content\_optimum.py}
records every survivor below $T=16288$.
There are seventeen before (G024.9), and exactly five
after it:
\[
\begin{array}{c|c|c|c|c|r}
 (a_1,a_2,g)&c&d&(c_0,c_1,M,h)&w&\text{relaxed score}\\\hline
 (6,9,3)&3&4&(2,1,4,0)&4&3716\\
 (2,1,1)&1&4&(6,3,4,1)&20&5180\\
 (2,1,1)&1&4&(6,7,4,0)&20&5180\\
 (10,5,5)&3&4&(2,1,4,0)&4&5268\\
 (6,15,3)&1&4&(2,1,20,0)&4&5372
\end{array}                                                         \tag{G024.10}
\]
No outer index has been enumerated or bounded.

\paragraph{Uniform exclusion of every survivor.}
All five rows have $d=4$. The effective parity-correct
primitive polynomial of that mass is unique:
\[
 \delta_0=q^2+q+q^{-1}+q^{-2}=D(D-3).
\]
Indeed these four odd coefficients already use the entire
mass. In the first, fourth and fifth rows $W_0=D$ and
$[q^0](D\delta_0)=2$. The first and fifth rows therefore
have $[q^0]G_1=12<16$, while the fourth has
$[q^0]G_2=10<14$. Each violates an actual trace target.

For the third row the coefficient list of $\delta_0^2$
is $(4,2,1,2,1)$ and $S=2W_0\delta_0^2$, so
$[q^0]S=2(6\cdot4+2\cdot7\cdot2)=104<144$.
For the second row $h=1$ means
\[
 W_0=6+3(q+q^{-1})+4(q^j+q^{-j})\quad(j\ge2)
\]
for an arbitrary integer $j$. The same list yields
\[
 [q^0]S=48+24+16[q^j]\delta_0^2\le104<144.
\]
Thus this entire unbounded-location family also fails.
All necessary-relaxation survivors below $5932$ have
been excluded. The actual table (G024.1) attains $5932$,
proving the asserted optimum.\hfill$\square$

\paragraph{Geometric translation and the remaining gate.}
NN-G010 realizes the integral effective multiplicities of
(G024.2) by actual graded sign representations in an
additive monoidal kernel category on a ten-point scheme.
The exact dictionaries remain SGA 3, I, 4.7.2--4.7.3
for graded representations and archive cote 2, scan 17
for finite-kernel composition. There is no object of
fractional dimension attached to $\Lambda_2$; the rational
factor is merely a coordinate factorization in $\mathbb Q(D)$.
The structure constant $\delta\Lambda_2$ is the actual
integral object. The content-three factor is three copies
of its effective cell object, not division by three in
an effective Grothendieck semiring.

For the equal-cell, constant-ratio tables realized by these
chosen kernel generators, the minimum sum over the $132$
coefficient positions of half the excess of dimension over
absolute signed trace is exactly $5932$. This is a minimum
for the specified character-valued multiplication data and
its NN-G010 realizations, not for all monoidal categories
or all generating sets with the same Grothendieck ring.

The same supplied trace algebra therefore admits smaller
composing effective lifts after a purely auxiliary
integrality restriction is removed. This demonstrates
why formal factor parameters cannot be identified with
canonical count objects. It establishes neither a
canonical family nor a rank-four basis, PH0, PH1 or
nonstandard purity. The precise remaining lower-score
gates are unequal general cells, nonconstant ratios,
or a different multiplication-table shape; arbitrary
common-cell degree and content are now covered.

